\begin{figure}
    \centering
    \begin{minipage}{0.49\textwidth}
        \centering
        \raisebox{0.1\height}{%
            \begin{tikzpicture}[scale=2]

  % Circle center M and radius R
  \coordinate (M) at (0,0);
  \def\R{0.8}

  % Point P and R
  \coordinate (P) at (1.2,0.8);
  \coordinate (R) at ({\R*cos(60)},{\R*sin(60)});

  % Inversion factor
  \pgfmathsetmacro{\k}{0.25}

  % Inverted point P'
  \coordinate (Pp) at (\k*1.2, \k*0.8); % same coordinates as P

  % Circle
  \draw[darkgreen, line width=1.8pt] (M) circle (\R);

  % Line through M, P', P
  \draw[black, thick] (M) -- (Pp) -- (P);
  \draw[black, thin] (M) -- (R);

  % Points
  \fill[blue] (P) circle (1pt);
  \fill[red] (Pp) circle (1pt);
  \fill[black] (M) circle (1pt);

  % Labels
  \node[left] at (M) {$M$};
  \node[below] at (P) {$P$};
  \node[below] at (Pp) {$P'$};
  \node[above right, yshift=-1pt, xshift=-1pt] at (R) {$R$};

\end{tikzpicture}
%
        }
    \end{minipage}
    \hfill
    \begin{minipage}{0.49\textwidth}
        \centering
        \begin{tikzpicture}[scale=2]

  % Circle center M and radius R
  \coordinate (M) at (0,0);
  \def\R{0.8}
  \def\r{0.3}

  % Circles and Line
  \draw[darkgreen, line width=1.8pt] (M) circle (\R);
  \draw[red, line width=1pt] ($(M)+(\r,0)$) circle (\r);
  \draw[blue, thick] ($(M)+(1.2,-1.2)$) -- ($(M)+(1.2,1.2)$);
  \draw[blue, thin, dashed] ($(M)+(1.2,-1.4)$) -- ($(M)+(1.2,1.4)$);

  % Points
  \fill[black] (M) circle (1pt);
  \fill[red] ({2*\r},0) circle (1pt) node[black, left] {$P'$};
  \fill[blue] (1.2,0) circle (1pt) node[black, right] {$P$};

  % Labels
  \node[left] at (M) {$M$};

\end{tikzpicture}

    \end{minipage}
    \caption{Illustration der Kreisspiegelung~\cite{kreisspiegelung} links
    mit einem roten Punkt $P'$ der in den blauen Punkt $P$ abgebildet wird. 
    Rechts ein Beispiel an einem roten Kreis, welcher in eine blaue Gerade abgebildet wird.
    Der grüne Kreis ist der sogenannte Inversionskreis dieser Transformation.}
    \label{fig:kreisspiegelung}
\end{figure}
